%% start of file `template.tex'.
%% Copyright 2006-2015 Xavier Danaux (xdanaux@gmail.com), 2020-2024 moderncv maintainers (github.com/moderncv).
%
% This work may be distributed and/or modified under the
% conditions of the LaTeX Project Public License version 1.3c,
% available at http://www.latex-project.org/lppl/.


\documentclass[11pt,a4paper,sans]{moderncv}        % possible options include font size ('10pt', '11pt' and '12pt'), paper size ('a4paper', 'letterpaper', 'a5paper', 'legalpaper', 'executivepaper' and 'landscape') and font family ('sans' and 'roman')

\moderncvcolor{black}
\moderncvstyle[]{banking}
\nopagenumbers{}                                  % uncomment to suppress automatic page numbering for CVs longer than one page

\usepackage{hyperref}

% adjust the page margins
\usepackage[hmargin=0.5in,vmargin=10pt]{geometry}

% font loading
% for luatex and xetex, do not use inputenc and fontenc
% see https://tex.stackexchange.com/a/496643
\ifxetexorluatex
  \usepackage{fontspec}
  \usepackage{unicode-math}
  \usepackage{FiraSans}
  \defaultfontfeatures{Ligatures=TeX}
  \setmainfont{Fira Sans}
  \setsansfont[ItalicFont={Fira Sans Light Italic},
            BoldFont={Fira Sans Bold},
            BoldItalicFont={Fira Sans Italic}]
            {Fira Sans Light}
  % \setmonofont{Latin Modern Mono}
  \setmonofont{Monaspace Neon}
  \setmathfont{Latin Modern Math}

  % you may also consider Fira Sans Light for a extra modern look
  %\setsansfont[ItalicFont={Fira Sans Light Italic},%
  %           BoldFont={Fira Sans},%
  %           BoldItalicFont={Fira Sans Italic}]%
  %           {Fira Sans Light}%
\else
  \usepackage[utf8]{inputenc}
  \usepackage[T1]{fontenc}
  \usepackage{lmodern}
\fi

% document language
\usepackage[UTF8,scheme=plain]{ctex}  % FIXME: using spanish breaks moderncv
\setCJKmainfont{SimHei}
\setCJKsansfont{SimHei}
% \setCJKmainfont{Noto Sans CJK SC}
% \setCJKsansfont{Noto Sans CJK SC}

\usepackage{setspace}
\usepackage{needspace}
\linespread{1.15}

% Compact headings and whitespace for moderncv's banking style.
\renewcommand*{\namefont}{\fontsize{24}{24}\selectfont\bfseries\upshape}
\renewcommand*{\sectionfont}{\fontsize{16}{14}\selectfont\bfseries\upshape}

\patchcmd{\makecvhead}{\vspace{2.5em}}{\vspace{1.5em}}
  {}{\PackageError{cv-layout}{Could not adjust title spacing}{}}
% This moderncv version defines section with xparse; patch its code macro.
\ExpandArgs{c}\patchcmd{section code}{\addvspace{2.5ex}}{\addvspace{2ex}}
  {}{\PackageError{cv-layout}{Could not adjust section spacing before}{}}
\ExpandArgs{c}\patchcmd{section code}{\addvspace{1ex}}{\addvspace{0.5ex}}
  {}{\PackageError{cv-layout}{Could not adjust section spacing after}{}}

% personal data
\name{\textmd{\textsc{吴胜桥}}}{}
\phone[mobile]{+86~159~8583~0350}
\email{Blind-Guess-Senior@outlook.com}
\homepage{https://blind-guess-senior.github.io/}

\social[github]{Blind-Guess-Senior}

% \extrainfo{additional information}                 % optional, remove / comment the line if not wanted
% \quote{Some quote}                                 % optional, remove / comment the line if not wanted

%----------------------------------------------------------------------------------
%            content
%----------------------------------------------------------------------------------
\begin{document}
% \begin{CJK*}{UTF8}{gbsn}                          % to typeset your resume in Chinese using CJK
%-----       resume       ---------------------------------------------------------
\makecvtitle

\section{教育经历}  % for 'contemporary' style use optional argument for displaying an icon, e.g. \section[\faGraduationCap]{Education}
\cventry{2024 -- 至今}{本科 | \textit{计算机科学与技术}}{上海科技大学}{\textit{中国 上海}}{}{}  % arguments 3 to 6 can be left empty

\vspace{0.5em}
在校期间已经修读操作系统、编译原理、计算机网络、程序分析等课程。

\needspace{8\baselineskip}
\section{项目经历}
\cventry{2026.10}{\textcolor[RGB]{128,128,128}{2026 聚光灯 GameJam 参赛项目} \textcolor[RGB]{84,174,255}{\href{https://github.com/Blind-Guess-Senior/Quququ}{Link}}}{Quququ \textmd{\textsc{| 团队项目}}}{Unity}{}{
  \begin{itemize}
    \item 设计并实现 AMP\&PHI 模块化角色架构及角色组装器，通过阶段 Hook 与结算单元支持插件组合和交互
    \item 拆分并跟进另外 2 名程序的开发任务，统筹接口与进度，团队完成完整作战流程及部分插件
  \end{itemize}}
% \vspace{0.1em}

\needspace{4\baselineskip}
\cventry{2026.07}{\textcolor[RGB]{128,128,128}{以“拍键盘”为核心输入的2D节奏游戏}  \textcolor[RGB]{84,174,255}{\href{https://github.com/Blind-Guess-Senior/KBre}{Link}}}{KBre \textmd{\textsc{| 团队项目}}}{Unity}{}{
  \begin{itemize}
    \item 玩家以“拍键盘”的方式控制角色移动、攻击。拍击、滑动、连续拍击等输入模式有各自的特殊效果。玩家需要配合游戏节拍输入以提高攻击伤害并降低资源消耗
    \item 设计整体代码框架，设计并实现事件系统、输入模式、场景物体、场景控制、敌人状态机等
  \end{itemize}}
% \vspace{0.1em}

\needspace{4\baselineskip}
\cventry{2026.03 -- 至今}{\textcolor[RGB]{128,128,128}{基于预编译对话脚本的 Unity 对话系统库} \textcolor[RGB]{84,174,255}{\href{https://github.com/Blind-Guess-Senior/ArtifactDialoguer}{Link}}}{Artifact Dialoguer \textmd{\textsc{| 个人项目}}}{C\#, .NET, Unity}{}{
  \begin{itemize}
    \item 设计简洁且无歧义的对话脚本语法，并实现 LL(1) 语法解析器将其编译为 Unity ScriptableObject
    \item 实现 Unity 内的运行时对话控制器，读取编译产物并驱动游戏内对话流程
  \end{itemize}
}
% \vspace{0.1em}

\needspace{4\baselineskip}
\cventry{2026.02 -- 至今}{\textcolor[RGB]{128,128,128}{轻量 Unity 工具库} \textcolor[RGB]{84,174,255}{\href{https://github.com/Blind-Guess-Senior/Unity-Utils}{Link}}}{Artifact Unity Utils \textmd{\textsc{| 个人项目}}}{C\#, .NET, Unity}{}{
  \begin{itemize}
    \item 基于 C\# 和 .NET 实现事件总线、依赖注入、有限状态机等通用架构与四叉树等数据结构
  \end{itemize}
}
% \vspace{0.1em}

\needspace{4\baselineskip}
\cventry{2025.07}{\textcolor[RGB]{128,128,128}{2D能力锁平台跳跃游戏} \textcolor[RGB]{84,174,255}{\href{https://github.com/Blind-Guess-Senior/C-17ES-CosmicHeart}{Link}}}{C-17ES-α \textmd{\textsc{| 团队项目}}}{Unity}{}{
  \begin{itemize}
    \item 玩家随流程逐步解锁四个方向的移动。不同方向的移动具有不同效果，玩家需要组合运用能力完成关卡
    \item 设计整体代码框架，设计并实现事件系统、移动能力、场景内交互物体等
  \end{itemize}}
% \vspace{0.1em}

% \needspace{8\baselineskip}
\section{工作经历}
\cventry{2026.05 -- 至今}{\textcolor[RGB]{128,128,128}{后端开发}}{影眸科技 \textmd{\textsc{| 实习}}}{\textit{上海}}{}{
  \begin{itemize}
    \item 接入多个广告事件回传平台，重构回传逻辑，提高准确性与灵活性
    \item 开发支付流程后端，实现订阅及支付相关新功能
  \end{itemize}}
% \vspace{0.1em}

% \needspace{3\baselineskip}
\section{比赛经历}
\cventry{2022.10}{CCF 非专业级软件能力认证提高组（CSP-S） 第二轮}{CCF CSP-S 第二轮 \textmd{\textsc{| 二等奖}}}{C++, Algorithm}{}{}

% \needspace{3\baselineskip}
\section{专业技能}
\cvitem{语言}{C\# | C/C++ | Javascript/Typescript | Nix | OCaml | Python | LaTeX}
\cvitem{引擎与框架}{Unity | .NET | React}

\section{个人情况}

热衷于游戏客户端工业化生产管线和高性能工具开发。热爱钻研操作系统和语言底层知识，并致力于以更本质的角度解决痛点问题。在多个项目中带队开发，积累了丰富的工程实践经验和项目管理经验。追求在任何场景下高效率地打造可信的生产力系统。

\clearpage

\end{document}


%% end of file `template.tex'.
